\documentclass[aspectratio=169, 12pt]{beamer}

% --- Theme ---
\usetheme{default}
\usecolortheme{orchid}
\usefonttheme{professionalfonts}

% Clean navigation
\setbeamertemplate{navigation symbols}{}
\setbeamertemplate{footline}[frame number]
\setbeamertemplate{itemize items}[circle]
\setbeamertemplate{enumerate items}[default]

% Colors
\definecolor{dblue}{RGB}{0,51,102}
\definecolor{mblue}{RGB}{0,85,150}
\definecolor{accent}{RGB}{180,40,40}
\setbeamercolor{structure}{fg=dblue}
\setbeamercolor{frametitle}{fg=dblue}
\setbeamercolor{title}{fg=dblue}
\setbeamercolor{block title}{bg=dblue!15, fg=dblue}
\setbeamercolor{block body}{bg=dblue!5}
\setbeamercolor{alerted text}{fg=accent}

% Packages
\usepackage{amsmath,amssymb,amsthm}
\usepackage{mathrsfs}
\usepackage{tikz}
\usetikzlibrary{arrows.meta, positioning, decorations.pathmorphing}
\usepackage{array}

% Theorem environments for beamer
\setbeamertemplate{theorems}[numbered]
\newtheorem{proposition}[theorem]{Proposition}
\newtheorem{claim}[theorem]{Claim}
\newtheorem{conjecture}[theorem]{Conjecture}

% Notation
\newcommand{\NN}{\mathcal{N}}
\newcommand{\CC}{\mathcal{C}}
\newcommand{\DD}{\mathcal{D}}
\newcommand{\CF}{\mathcal{C}_F}
\newcommand{\vchi}{\vec{\chi}_B}
\DeclareMathOperator{\dom}{dom}

% Title
\title{The Borel dichromatic number}
\author{Tonatiuh Matos-Wiederhold}
\institute{University of Toronto}
\date{Caltech Logic Seminar \\ April 29, 2026}

\begin{document}

% =====================================================================
\begin{frame}
\titlepage
\end{frame}
% =====================================================================

% =====================================================================
% PART I: BACKGROUND
% =====================================================================

\begin{frame}{Outline}
\begin{enumerate}
	\item \textbf{Background:} Borel chromatic numbers and basis problems
	\item \textbf{Dichromatic number:} a directed analogue
	\item \textbf{Classical reduction:} from hypergraphs to digraphs
	\item \textbf{Borel lifting:} Thornton's coding framework
	\item \textbf{Open questions}
\end{enumerate}
\end{frame}

% --- Borel graphs ---

\begin{frame}{Borel graphs}
A \alert{Borel graph} $G = (X, E)$ consists of:
\begin{itemize}
	\item a standard Borel space $X$ (vertex set),
	\item a Borel symmetric irreflexive relation $E \subseteq X^2$ (edges).
\end{itemize}

\medskip

A \alert{Borel $k$-coloring} is a Borel map $c\colon X \to k$ with $c(x) \neq c(y)$ for all $(x,y) \in E$.

\medskip

The \alert{Borel chromatic number} $\chi_B(G)$ is the minimum $k$ admitting a Borel $k$-coloring.

\medskip
\begin{block}{Basic fact (folklore)}
Every Borel graph of degree $\leq d$ admits a Borel $(d+1)$-coloring.
\end{block}
\end{frame}

% --- Why Borel? ---

\begin{frame}{Why Borel?}
Every Borel graph has a (possibly non-Borel) proper coloring.

\medskip

The point: \alert{requiring the coloring to be Borel} makes the problem genuinely harder.

\medskip

\begin{example}
There exist Borel graphs with $\chi(G)=2$ but $\chi_B(G) > \aleph_0$.
\end{example}

\medskip

The Borel chromatic number captures a notion of \emph{definable complexity} of the coloring problem.
\end{frame}

% --- The G0 dichotomy ---

\begin{frame}{The $\mathbb{G}_0$-dichotomy (Kechris--Solecki--Todorčević, 1999)}

\begin{theorem}[KST '99]
There exists a Borel graph $\mathbb{G}_0$ such that, for every analytic graph $G$:
$$
\chi_B(G) > \aleph_0
\quad\Longleftrightarrow\quad
\mathbb{G}_0 \leq_B G
$$
where $\leq_B$ denotes the existence of a Borel homomorphism.
\end{theorem}

\pause
\medskip

In other words: $\{\mathbb{G}_0\}$ is a \alert{one-element basis} for the class of Borel graphs with uncountable Borel chromatic number.

\medskip

\textbf{Interpretation:} To prove $\chi_B(G) > \aleph_0$, there is essentially \emph{one} obstruction to find. The problem is ``well-structured.''
\end{frame}

% --- What is a basis? ---

\begin{frame}{What is a basis?}
Fix a class $\mathscr{C}$ of Borel graphs (e.g.\ those with $\chi_B \geq k$).

\medskip

A family $\mathbf{B}$ is a \alert{basis} for $\mathscr{C}$ under Borel homomorphism if:
$$
G \in \mathscr{C} \quad\Longleftrightarrow\quad \exists\, H \in \mathbf{B},\ H \leq_B G.
$$

\pause
\medskip

\textbf{Small basis = ``simple'' decision problem:}
\begin{itemize}
	\item Size 1: one obstruction to check (like $\mathbb{G}_0$)
	\item Finite: finitely many obstructions
	\item Countable: ``tame'' complexity
	\item Uncountable or none: the problem is genuinely complex
\end{itemize}
\end{frame}

% --- The L0 dichotomy ---

\begin{frame}{The $\mathbb{L}_0$-dichotomy}
\begin{theorem}[Carroy--Miller--Schrittesser--Vidnyánszky]
There exists a Borel graph $\mathbb{L}_0$ such that, for every analytic graph $G$:
$$
\chi_B(G) > 2
\quad\Longleftrightarrow\quad
\mathbb{L}_0 \leq_B G.
$$
\end{theorem}

Again a one-element basis, this time for the $2$ vs.\ $3$-color threshold.

\pause
\medskip

\textbf{So far:} at both the $2/3$-color threshold and the countable/uncountable threshold, there is a dichotomy with a one-element basis.

\medskip

\textbf{Question:} Does this continue for higher thresholds?
\end{frame}

% --- TV: no more dichotomies ---

\begin{frame}{No dichotomy at higher thresholds}
\begin{theorem}[Todorčević--Vidnyánszky, 2021]
For every $k \geq 3$, the set of (codes for) Borel graphs admitting a Borel $k$-coloring is $\mathbf{\Sigma}^1_2$-complete.
\end{theorem}

\pause
\medskip

\begin{corollary}
For $k \geq 3$: no countable basis for the class of Borel graphs with $\chi_B(G) > k$.
\end{corollary}

\medskip

The landscape for \emph{undirected} Borel chromatic numbers is completely understood:
\begin{center}
\begin{tabular}{ll}
	$\chi_B > 2$ & basis of size $1$ \\
	$\chi_B > k$, $k\geq 3$ & no countable basis \\
	$\chi_B > \aleph_0$ & basis of size $1$ \\
\end{tabular}
\end{center}
\end{frame}

% --- Sigma12 completeness, informally ---

\begin{frame}{Interlude: what does $\mathbf\Sigma^1_2$-completeness mean?}
Recall the \alert{projective hierarchy}:

\medskip

\begin{center}
\begin{tikzpicture}[scale=0.9]
	\node (s11) at (0,0) {$\Sigma^1_1$ (analytic)};
	\node (p11) at (4,0) {$\Pi^1_1$ (coanalytic)};
	\node (d11) at (2,-0.8) {$\Delta^1_1$ (Borel)};
	\node (s12) at (0,2) {$\Sigma^1_2$};
	\node (p12) at (4,2) {$\Pi^1_2$};
	\node (d12) at (2,1.2) {$\Delta^1_2$};
	\draw[->] (d11) -- (s11);
	\draw[->] (d11) -- (p11);
	\draw[->] (s11) -- (d12);
	\draw[->] (p11) -- (d12);
	\draw[->] (d12) -- (s12);
	\draw[->] (d12) -- (p12);
\end{tikzpicture}
\end{center}

\medskip

A set $A \subseteq \omega$ is \alert{$\Sigma^1_2$} if there is a $\Pi^1_1$ set $P \subseteq \omega \times \NN$ with:
$$
n \in A \quad\Longleftrightarrow\quad \exists x \in \NN,\ (n,x) \in P.
$$

\end{frame}

\begin{frame}{$\Sigma^1_2$-completeness: the key example}
\begin{block}{Why ``admits a Borel $k$-coloring'' is $\Sigma^1_2$}
Given a \emph{nice code} $e_G$ for a Borel graph $G$:
$$
\text{``$G$ admits a Borel $k$-coloring''}
\quad\Longleftrightarrow\quad
\exists\, e_c \in \omega\; \bigl[\underbrace{\text{``$\mathcal{D}_{e_c}$ is a $k$-coloring of $G$''}}_{\Pi^1_1 \text{ in } e_c, e_G}\bigr]
$$
\end{block}

\pause
\medskip

The inner condition is $\Pi^1_1$ because it asserts:
\begin{itemize}
	\item totality: $\forall x \in V(G),\ \exists! i < k,\ c(x) = i$ \quad [$\Pi^1_1$]
	\item properness: $\forall (x,y) \in E(G),\ c(x) \neq c(y)$ \quad [$\Pi^1_1$]
\end{itemize}

\medskip

The outer $\exists e_c$ quantifier makes the whole thing $\Sigma^1_2$.

\pause
\medskip

\alert{$\Sigma^1_2$-complete} means: this set is \emph{as hard as any} $\Sigma^1_2$ set. Every $\Sigma^1_2$ set reduces to it via a $\Delta^1_1$ function.
\end{frame}

\begin{frame}{Why does $\Sigma^1_2$-completeness destroy countable bases?}
\begin{proof}[Proof sketch (for chromatic number)]
Suppose $\mathbf{B} = \{H_1, H_2, \ldots\}$ is a countable basis for $\chi_B(G) > k$. Then:
$$
\chi_B(G) > k
\quad\Longleftrightarrow\quad
\exists\, n,\ \exists\text{ Borel homomorphism } H_n \to G.
$$

\medskip

For each fixed $H_n$, ``$\exists$ Borel hom $H_n \to G$'' is $\Sigma^1_2$ in the code for $G$.

\medskip

A countable union of $\Sigma^1_2$ sets is $\Sigma^1_2$.

\medskip

But the left-hand side $\{\text{code}(G) : \chi_B(G) > k\}$ is $\Pi^1_2$-complete.

\medskip

This gives $\Pi^1_2 \subseteq \Sigma^1_2$, contradicting the strictness of the projective hierarchy.
\end{proof}
\end{frame}

% =====================================================================
% PART II: DICHROMATIC NUMBER
% =====================================================================

\begin{frame}{The dichromatic number (Neumann-Lara, 1982)}
Let $D = (V, A)$ be a directed graph.

\medskip

A \alert{dicoloring} is a map $c\colon V \to k$ such that every color class induces an \emph{acyclic} subdigraph (no directed cycles).

\medskip

The \alert{dichromatic number} $\vec\chi(D)$ is the minimum such $k$.

\pause
\medskip

\begin{example}
\begin{itemize}
	\item A tournament on 3 vertices forming a directed 3-cycle: $\vec\chi = 2$.
	\item A directed graph with no directed cycles: $\vec\chi = 1$.
	\item If $D$ is a symmetric digraph ($\{u,v\} \in E \Leftrightarrow (u,v),(v,u) \in A$): $\vec\chi(D) = \chi(\text{underlying graph})$.
\end{itemize}
\end{example}
\end{frame}

\begin{frame}{Dichromatic vs.\ chromatic: key difference}
\begin{block}{Classical complexity}
\begin{itemize}
	\item ``Is $\chi(G) \leq 2$?'' is decidable in polynomial time (bipartiteness).
	\item ``Is $\vec\chi(D) \leq 2$?'' is \alert{NP-complete} \quad (Bokal et al., 2004).
\end{itemize}
\end{block}

\medskip

The dichromatic number is ``harder'' than the chromatic number: complexity starts already at the $2/3$-color threshold.

\pause
\medskip

\textbf{Key observation:} Acyclicity is \emph{not} a local constraint.

Checking that a color class is independent (for chromatic number) requires looking at edges one at a time.

Checking that a color class is acyclic may require following long directed paths.

\medskip

In particular, dicoloring is \alert{not a CSP} since acyclicity cannot be expressed as a homomorphism to a fixed finite template.
\end{frame}

\begin{frame}{Borel dichromatic number}
For Borel directed graphs, define $\vchi(D)$ analogously, requiring the dicoloring to be Borel.

\medskip

\begin{theorem}[Raghavan--Xiao, 2024]
There exists a family $\{H_\alpha\}_{\alpha < \mathfrak{c}}$ of Borel digraphs, of size continuum, such that:
$$
\vchi(D) > \aleph_0
\quad\Longleftrightarrow\quad
\exists\, \alpha,\ H_\alpha \leq_B D.
$$
The $H_\alpha$ are pairwise Borel-incomparable.
\end{theorem}

\pause
\medskip

This is the directed analogue of the $\mathbb{G}_0$-dichotomy, but the basis has \emph{continuum} many elements instead of one.

\medskip

\textbf{Open:} Is $\mathfrak{c}$ optimal, or could a smaller (perhaps countable) basis exist?
\end{frame}

\begin{frame}{The main result}
\begin{proposition}
The set of nice codes for Borel directed graphs admitting a Borel $2$-dicoloring is $\mathbf\Sigma^1_2$-complete.

Equivalently, the set of nice codes for Borel directed graphs with $\vchi(D) \geq 3$ is $\mathbf\Pi^1_2$-complete.
\end{proposition}

\pause

\begin{corollary}
There is no countable basis for Borel digraphs of Borel dichromatic number $\geq 3$.
\end{corollary}

\pause

\begin{block}{Recovering the chromatic number}
Given an undirected graph $G$, symmetrize it: $\vec{G}$ has arcs $(u,v)$ and $(v,u)$ for each edge $\{u,v\}$. Then $\vchi(\vec{G}) = \chi_B(G)$.

This reduction is $\Delta^1_1$ on codes, so $\Sigma^1_2$-completeness of $k$-dicoloring for $k \geq 4$ follows from Todorčević--Vidnyánszky.
\end{block}
\end{frame}

\begin{frame}{The complete picture}
\begin{center}
\renewcommand{\arraystretch}{1.3}
\begin{tabular}{lll}
	\hline
	\textbf{Class} & \textbf{Basis?} & \textbf{Reference} \\
	\hline
	$\chi_B(G) > \aleph_0$ & Yes, size 1 ($\mathbb{G}_0$) & KST '99 \\
	$\chi_B(G) > 2$ & Yes, size 1 ($\mathbb{L}_0$) & CMSV \\
	$\chi_B(G) > k$, $k\geq 3$ & \alert{No} ($\mathbf\Pi^1_2$-complete) & TV '21 \\
	\hline
	$\vchi(D) > \aleph_0$ & Yes, size $\mathfrak{c}$ & RX '24 \\
	$\vchi(D) > k$, $k\geq 2$ & \alert{No} ($\mathbf\Pi^1_2$-complete) & \textbf{this talk} \\
	\hline
\end{tabular}
\end{center}

\pause
\medskip

The case $k=2$ for directed graphs is the \alert{last missing piece}.

\medskip

The proof: lift the classical NP-completeness reduction to the Borel setting.
\end{frame}

% =====================================================================
% PART III: THE CLASSICAL REDUCTION
% =====================================================================

\begin{frame}{Strategy}
\begin{center}
\begin{tikzpicture}[
	box/.style={draw, rounded corners, minimum width=3.5cm, minimum height=1cm, align=center, font=\small},
	arr/.style={-Stealth, thick}
]
	\node[box, fill=dblue!10] (hyp) at (0,0) {3-uniform hypergraph\\2-coloring};
	\node[box, fill=accent!10] (dig) at (7,0) {Digraph\\2-dicoloring};
	\node[box, fill=dblue!10] (bhyp) at (0,-2.5) {Borel 3-uniform\\hypergraph 2-coloring};
	\node[box, fill=accent!10] (bdig) at (7,-2.5) {Borel digraph\\2-dicoloring};

	\draw[arr] (hyp) -- node[above, font=\footnotesize] {Bokal et al.} (dig);
	\draw[arr] (bhyp) -- node[above, font=\footnotesize] {this talk} (bdig);
	\draw[arr, dashed, gray] (hyp) -- node[left, font=\footnotesize, gray] {} (bhyp);

	\node[font=\footnotesize, text=mblue] at (0,-3.5) {$\Sigma^1_2$-complete (Thornton)};
	\node[font=\footnotesize, text=accent] at (7,-3.5) {$\Sigma^1_2$-complete (\textbf{new})};
\end{tikzpicture}
\end{center}
\end{frame}

\begin{frame}{Step 1: The classical starting point}
\begin{fact}
The problem of deciding whether a 3-uniform hypergraph admits a 2-coloring is NP-complete.
\end{fact}

\medskip

A \alert{3-uniform hypergraph} $H = (V, E)$ has hyperedges $e \in \binom{V}{3}$.

A \alert{2-coloring} is a map $c\colon V \to 2$ such that no hyperedge is monochromatic.

\pause
\medskip

\textbf{Goal:} Reduce this to 2-dicoloring of directed graphs.

\medskip

\textbf{Idea:} Build a directed graph $D(H)$ from $H$ using a finite \emph{gadget} that translates ``no monochromatic hyperedge'' into ``no monochromatic directed cycle.''
\end{frame}

\begin{frame}{The template gadget}
\begin{columns}
\begin{column}{0.45\textwidth}
The gadget $F$ has:
\begin{itemize}
	\item \alert{Shared vertices} $\{a,b,c\}$ (will become hypergraph vertices)
	\item \alert{Auxiliary vertices} $\{a',b',c'\}$ (local to each gadget copy)
\end{itemize}

\medskip

\begin{lemma}[Gadget property]
A map $\sigma\colon\{a,b,c\}\to 2$ extends to a 2-dicoloring of~$F$ iff $\sigma$ is \alert{not constant}.
\end{lemma}
\end{column}
\begin{column}{0.5\textwidth}
	\begin{center}
		\input{figures/bokal-gadget.tex}
	\end{center}
\end{column}
\end{columns}

\pause
\medskip

\begin{center}
\small
\begin{tabular}{c|c}
	coloring of $\{a,b,c\}$ & extends to $F$? \\
	\hline
	$(0,0,1)$, $(0,1,0)$, $(1,0,0)$, \ldots & \alert{Yes} \\
	$(0,0,0)$ or $(1,1,1)$ & \alert{No}
\end{tabular}
\end{center}
\end{frame}

\begin{frame}{The reduction $H \mapsto D(H)$}
Fix a 3-uniform hypergraph $H = (V, E)$ and a linear order on $V$.

\medskip

\textbf{Sorted hyperedges:} $E_{\mathrm{sort}} = \{(v_1,v_2,v_3) \in E : v_1 < v_2 < v_3\}$.

\medskip

\textbf{Vertex set of $D(H)$:}
$$
V_D = V \;\cup\; \bigl\{(e, t) : e \in E_{\mathrm{sort}},\; t \in \{a', b', c'\}\bigr\}
$$

\textbf{Arc set:} For each $e = (v_1,v_2,v_3)$, place a copy $F_e$ of the gadget:
\begin{itemize}
	\item shared labels $a, b, c \mapsto v_1, v_2, v_3$
	\item auxiliary labels $a', b', c' \mapsto (e, a'), (e, b'), (e, c')$
\end{itemize}

\medskip

\alert{Key:} Shared vertices are identified across gadgets. Auxiliary vertices are disjoint.
\end{frame}

\begin{frame}{Forward direction: $H$ 2-colorable $\Rightarrow$ $D(H)$ 2-dicolorable}
Let $c\colon V \to 2$ be a 2-coloring of $H$.

\medskip

For each hyperedge $e$: the shared vertices get a non-constant coloring, so by the gadget property it \alert{extends} to a 2-dicoloring of $F_e$.

\medskip

Auxiliary vertices of different gadgets are disjoint $\Rightarrow$ the extensions combine into a global 2-dicoloring of $D(H)$.

\pause
\bigskip

\begin{center}
\begin{tikzpicture}
	\node[draw, circle, fill=blue!20, minimum size=0.7cm] (v1) at (0,0) {$\color{accent}0$};
	\node[draw, circle, fill=blue!20, minimum size=0.7cm] (v2) at (2,0) {$\color{accent}0$};
	\node[draw, circle, fill=blue!20, minimum size=0.7cm] (v3) at (4,0) {$\color{accent}1$};
	\node[font=\small] at (2,-0.7) {non-constant $\checkmark$};
	\node[font=\small] at (6,0) {$\Rightarrow$ extends to $F_e$};
\end{tikzpicture}
\end{center}
\end{frame}

\begin{frame}{Reverse direction: $D(H)$ 2-dicolorable $\Rightarrow$ $H$ 2-colorable}
Let $\psi\colon V_D \to 2$ be a 2-dicoloring of $D(H)$.

\medskip

Define $\varphi = \psi|_V$ (just \alert{restrict} to the original vertices).

\medskip

For each hyperedge $e = (v_1,v_2,v_3)$:
\begin{itemize}
	\item $F_e$ is a subdigraph of $D(H)$
	\item So $\psi|_{V(F_e)}$ is a valid 2-dicoloring of $F_e$
	\item By the gadget property: $\varphi(v_1), \varphi(v_2), \varphi(v_3)$ are \alert{not all equal}
\end{itemize}

\medskip

Hence $\varphi$ is a 2-coloring of $H$. \hfill$\square$

\pause
\bigskip

\begin{alertblock}{Important observation}
The reverse direction is a \emph{plain restriction}. No choice, no selection, no local finiteness needed. This will matter in the Borel setting.
\end{alertblock}
\end{frame}

% =====================================================================
% PART IV: THE BOREL LIFTING
% =====================================================================

\begin{frame}{Borel lifting: the goal}
\begin{fact}[Thornton]
The set of nice codes for locally finite Borel 3-uniform hypergraphs admitting a Borel 2-coloring is $\mathbf\Sigma^1_2$-complete.
\end{fact}

\pause
\medskip

\textbf{We need to verify:}
\begin{enumerate}
	\item The map $\mathrm{code}(H) \mapsto \mathrm{code}(D(H))$ is $\Delta^1_1$.
	\item ``$H$ is Borel 2-colorable'' $\Longleftrightarrow$ ``$D(H)$ is Borel 2-dicolorable'' with \alert{Borel witnesses}.
\end{enumerate}

\pause
\medskip

\begin{block}{Why ``$\Delta^1_1$ on codes''?}
A $\Delta^1_1$ reduction from a $\Sigma^1_2$-complete set preserves $\Sigma^1_2$-hardness.

So if the reduction is $\Delta^1_1$ and the equivalence holds, the target is also $\Sigma^1_2$-hard.
\end{block}
\end{frame}

\begin{frame}{Thornton's coding framework (brief)}
Fix a good $\omega$-parametrization $U \subseteq \omega \times \NN$ of $\Pi^1_1$.

\medskip

A \alert{nice code} for a $\Delta^1_1$ set $B$ is an index $e$ such that $B$ is simultaneously described by a $\Pi^1_1$ formula and a $\Sigma^1_1$ formula, and these descriptions agree.

\medskip

\textbf{The key tool:}

\begin{lemma}[Toolkit]
The following operations are $\Delta^1_1$ in the codes:
\begin{columns}[t]
\begin{column}{0.45\textwidth}
\begin{enumerate}
	\item $(A,B) \mapsto A \cup B$
	\item $(A,B) \mapsto A \cap B$
	\item $(A,B) \mapsto A \times B$
\end{enumerate}
\end{column}
\begin{column}{0.53\textwidth}
\begin{enumerate}\setcounter{enumi}{3}
	\item $A \mapsto \dom(A)$
	\item $(A,f) \mapsto f(A)$, $f$ ctble-to-one
	\item $A \mapsto$ enum.\ of finite sections
\end{enumerate}
\end{column}
\end{columns}
\end{lemma}

\medskip

\textbf{Recipe:} Express the construction using these operations $\Rightarrow$ it is $\Delta^1_1$ on codes.
\end{frame}

\begin{frame}{The construction is $\Delta^1_1$ on codes}
Let $H$ be a locally finite Borel 3-uniform hypergraph, coded by $(V, E)$. Fix a $\Delta^1_1$ linear order $\preceq$ on $\NN$.

\bigskip

\textbf{Step 1: Sort hyperedges.}
$$E_{\mathrm{sort}} = \{(a,b,c) \in E : a \preceq b \preceq c\}$$
\hfill intersection with graph of $\preceq$ $\Rightarrow$ toolkit item (2) \checkmark

\pause
\bigskip

\textbf{Step 2: Build vertex set.}
$$V_D = V \cup (E_{\mathrm{sort}} \times T_{\mathrm{aux}})$$
\hfill product + union $\Rightarrow$ toolkit items (1), (3) \checkmark

\pause
\bigskip

\textbf{Step 3: Build arc set.} (next slide)
\end{frame}

\begin{frame}{Building the arc set}
Define the \alert{substitution map} $\varphi_H\colon E_{\mathrm{sort}} \times E_T \to V_D^2$:
$$
(e, (u,v)) \;\longmapsto\; \text{(corresponding arc of $D(H)$)}
$$
\begin{itemize}
	\item shared label in position $i$ $\mapsto$ vertex $v_i$ from $e = (v_1,v_2,v_3)$ \quad (projection)
	\item auxiliary label $\ell$ $\mapsto$ $(e, \ell)$ \quad (pairing)
\end{itemize}

\pause
\medskip

$\varphi_H$ is $\Delta^1_1$ in the code of $H$: it uses projections and products, which are recursive.

\medskip

$\varphi_H$ is \alert{finite-to-one}: $|E_T|$ is fixed, and each vertex appears in finitely many hyperedges (local finiteness).

\medskip

Arc set: $A_D = \varphi_H(E_{\mathrm{sort}} \times E_T)$ \hfill $\Rightarrow$ toolkit item (5) \checkmark

\pause
\bigskip

\begin{lemma}
The map $\mathrm{code}(H) \mapsto \mathrm{code}(D(H))$ is $\Delta^1_1$.
\end{lemma}
\end{frame}

\begin{frame}{Borel forward direction: Luzin--Novikov}
\begin{claim}
If $H$ is Borel 2-colorable, then $D(H)$ is Borel 2-dicolorable.
\end{claim}

\pause

\begin{proof}
Let $c\colon V \to 2$ be a Borel 2-coloring of $H$.
For each $e \in E_{\mathrm{sort}}$, at least one extension to a 2-dicoloring of $F_e$ exists.

\medskip

Define the set of valid local completions:
$$
S = \{(e, \sigma) \in E_{\mathrm{sort}} \times \CF : \sigma \text{ extends to a 2-dicoloring of } F_e\}
$$

$S$ is Borel, with \alert{finite} nonempty sections $S_e$.
By the \alert{Luzin--Novikov theorem}: $\exists$~Borel selector $e \mapsto \sigma(e) \in S_e$.

\medskip

Pasting these selections with $c$ on shared vertices gives a Borel 2-dicoloring of $D(H)$.
\end{proof}
\end{frame}

\begin{frame}{Borel reverse direction: just restrict}
\begin{claim}
If $D(H)$ is Borel 2-dicolorable, then $H$ is Borel 2-colorable.
\end{claim}

\pause

\begin{proof}
Let $\psi\colon V_D \to 2$ be a Borel 2-dicoloring of $D(H)$.

\medskip

Define $\varphi = \psi|_V$.

\medskip

Since $V \subseteq V_D$ is Borel, $\varphi$ is Borel.

\medskip

For each hyperedge $e$: $F_e \subseteq D(H)$ is a subdigraph, so $\psi|_{V(F_e)}$ is a valid 2-dicoloring, hence the shared-vertex colors are not all equal.

\medskip

So $\varphi$ is a Borel 2-coloring of $H$.
\end{proof}
\end{frame}

\begin{frame}{Proof of the main result}
\begin{proof}[Proof of Proposition]

\textbf{Hardness:} We have a $\Delta^1_1$ reduction from Borel 2-colorable 3-uniform hypergraphs ($\Sigma^1_2$-complete) to Borel 2-dicolorable digraphs. Hence the latter is $\Sigma^1_2$-hard.

\pause
\bigskip

\textbf{Upper bound:}
$$
\text{``$D$ admits a Borel 2-dicoloring''}
\quad\Longleftrightarrow\quad
\exists\, e_c\; [\underbrace{\text{``$\mathcal{D}_{e_c}$ is a valid 2-dicoloring''}}_{\Pi^1_1}]
$$

The $\Pi^1_1$ condition checks:
\begin{itemize}
	\item totality and range $\subseteq \{0,1\}$ \quad ($\Pi^1_1$)
	\item each color class induces no directed cycle \quad ($\Pi^1_1$: universal over finite paths)
\end{itemize}

\medskip

So the set is $\Sigma^1_2$. Combined with hardness: \alert{$\Sigma^1_2$-complete}. \hfill$\square$
\end{proof}
\end{frame}

\begin{frame}{The corollary: no countable basis}
\begin{proof}[Proof of Corollary]
Suppose $\mathbf{B} = \{H_1, H_2, \ldots\}$ is a countable basis for $\vchi(D) \geq 3$:
$$
\vchi(D) \geq 3
\quad\Longleftrightarrow\quad
\exists\, n,\ \exists\text{ Borel hom } H_n \to D.
$$

\pause
\medskip

For fixed $H_n$: ``$\exists$ Borel hom $H_n \to D$'' is $\Sigma^1_2$ in the code of $D$.

\medskip

Countable union $\Rightarrow$ the right-hand side is $\Sigma^1_2$.

\medskip

But the left-hand side is $\Pi^1_2$-complete (complement of our $\Sigma^1_2$-complete set).

\medskip

$\Pi^1_2 \subseteq \Sigma^1_2$ contradicts the strictness of the projective hierarchy. \hfill$\square$
\end{proof}
\end{frame}

\begin{frame}{A stronger conclusion}
The complexity argument rules out more than just a countable basis.

\medskip

\begin{block}{Observation}
If $\mathbf{B} \in \Sigma^1_2$ as a subset of the code space, then:
\begin{itemize}
	\item ``$H \in \mathbf{B}$'' is $\Sigma^1_2$
	\item ``$\exists$ Borel hom $H \to D$'' is $\Sigma^1_2$
	\item conjunction and $\exists$-quantification preserve $\Sigma^1_2$
\end{itemize}
So the right-hand side would be $\Sigma^1_2$, contradicting $\Pi^1_2$-completeness.
\end{block}

\pause
\medskip

\textbf{Conclusion:} Any basis $\mathbf{B}$ must satisfy $\mathbf{B} \notin \Sigma^1_2$. It is at least as complex to describe as the set of all digraphs without a Borel acyclic 2-coloring.

\medskip

\alert{A complexity lower bound destroys all ``simple'' bases, not just countable ones.}
\end{frame}

% =====================================================================
% PART V: OPEN QUESTIONS
% =====================================================================

\begin{frame}{Open questions I: the uncountable threshold}
Raghavan--Xiao give a continuum-size basis for $\vchi(D) > \aleph_0$.

\medskip

\begin{enumerate}
	\item[\textbf{Q1.}] \textbf{Is the basis irredundant?} I.e., does removing any single element break the basis? (Likely yes, from pairwise incomparability.)

	\bigskip

	\item[\textbf{Q2.}] \textbf{Can the basis be made smaller?} Is $\mathfrak{c}$ optimal?

	\bigskip

	\item[\textbf{Q3.}] \textbf{Is there a direct (non-complexity) proof} that no \emph{countable} basis exists for $\vchi(D) > \aleph_0$? Perhaps via diagonalization, adapting Miller's anti-basis argument for $G_f$ graphs.
\end{enumerate}

\medskip

\small
Note: the complexity approach does \emph{not} work here---the continuum-size basis shows the set is not $\Pi^1_2$-complete.
\end{frame}

\begin{frame}{Open questions II: structure and CSPs}
\begin{enumerate}\setcounter{enumi}{3}
	\item[\textbf{Q4.}] \textbf{Structural restrictions.} Does a countable basis exist for digraphs $G_f$ (generated by a single Borel function) with $\vec\chi_B(G_f) \geq 3$?

	Miller showed no countable basis exists for undirected $G_f$ with $\chi_B \geq 3$.

	\bigskip

	\item[\textbf{Q5.}] \textbf{Bounded width CSPs}. Thornton showed that width 1 structures are exactly those where solvable Borel instances have Borel solutions. Grebík--Vidnyánszky showed the Borel/classical CSP dichotomies diverge.

	\emph{What is the exact complexity dividing line within bounded width?}

	\bigskip

	\item[\textbf{Q6.}] \textbf{Measurable arboricity} (Thornton). Arboricity = undirected analogue of dichromatic number.

	\emph{Can $\mu$-measurable arboricity differ from classical arboricity by more than 1?}
\end{enumerate}
\end{frame}

\begin{frame}{}
\begin{center}
	\Huge Thank you.

	\bigskip
	\large
	\texttt{arXiv:} \textit{\href{https://arxiv.org/abs/2604.05228}{arxiv.org/abs/2604.05228}}
\end{center}
\end{frame}

\end{document}
